% Generated draft: requires mathematical and source review.
\documentclass[11pt]{article}
\usepackage[margin=25mm]{geometry}
\usepackage{fontspec}
\setmainfont{Latin Modern Roman}
\usepackage{amsmath,amssymb,mathtools,unicode-math}
\setmathfont{Latin Modern Math}
\usepackage{enumitem,needspace,xurl,hyperref}
\setlength{\parindent}{0pt}
\setlength{\parskip}{5pt}
\setlength{\emergencystretch}{3em}
\newcommand{\N}{\mathbb{N}}
\newcommand{\Z}{\mathbb{Z}}
\newcommand{\Q}{\mathbb{Q}}
\newcommand{\R}{\mathbb{R}}
\newcommand{\C}{\mathbb{C}}
\newcommand{\F}{\mathbb{F}}
\newcommand{\A}{\mathbb{A}}
\newcommand{\PP}{\mathbb{P}}
\newcommand{\E}{\mathbb{E}}
\newcommand{\eps}{\varepsilon}
\newcommand{\dd}{\,\mathrm{d}}
\begin{document}
\section*{UnsolvedMath ID: 30006376}
\textbf{Source problem code:} catalogue entry
\subsection{Definitions and mathematical statement}
\textbf{Unrecoverable statement.} The supplied record for permanent UnsolvedMath ID \(30006376\), source code \texttt{catalogue entry}, is titled ``Truncated Algebraic-K-Theory Abstract Fragment.'' Its extraction begins with the isolated word ``conjecture.'' and continues with workshop publicity about advances and applications of algebraic \(K\)-theory to neighboring areas of mathematics. No conjecture, mathematical objects, hypotheses, domains, boundary conditions, quantifiers, or conclusion are identifiable. Accordingly, no precise mathematical statement or associated definitions can be recovered from the supplied evidence.

The supplied statement status is \texttt{unrecoverable}, and no clean statement is supplied. The registry's verification note reports that the official report identifies the extraction as the tail of a sentence from the workshop abstract, rather than a standalone conjecture; this is a supplied verification claim, not an independent verification here. The registry labels the entry ``open,'' but that classification does not establish an identifiable open mathematical proposition. The intended conjecture remains unspecified.
\subsection{Short English statement}
This entry preserves a truncated workshop-abstract fragment concerning algebraic \(K\)-theory. The supplied text does not identify a mathematical claim, so a precise problem statement requires recovery and human review of the intended source passage. The supplied ``open'' label is registry metadata.
\subsection{Sources}
\begin{itemize}
\item Supplied local UnsolvedMath registry record, permanent ID \(30006376\), source code \texttt{catalogue entry}, title ``Truncated Algebraic-K-Theory Abstract Fragment,'' in the ``Oberwolfach Reports Open Problems'' collection. No proposer is identified; the supplied proposed year is \(2025\).
\item Supplied source citation: \emph{Algebraic K-theory (2026)}, Oberwolfach Reports, DOI \texttt{10.4171/owr/2025/36}, pp.\ 1949--2012. Supplied source URL and literature-source link labeled ``Oberwolfach report'': \url{https://doi.org/10.4171/owr/2025/36}. This report reference supplies provenance, but the supplied fragment does not identify a standalone conjecture or its location within those pages.
\item Literature assessment dated exactly \texttt{2026-08-22}: the supplied curation labels the entry ``open'' and states that the truncated occurrence of ``conjecture'' contains no identifiable mathematical claim and cannot be verified beyond that classification. This is supplied dated curation, not independently verified literature assessment; its background explicitly calls for independent verification of time-sensitive status and cited claims.
\item The supplied registry verification note says that the extraction comes from the tail of a workshop-abstract sentence, that no separate clean statement is promoted, and that the original record remains publicly visible and retained for provenance. These are supplied source-verification statements, not claims of source access here.
\item The extracted text supplies Mathematics Subject Classification (2020) codes \texttt{14-XX} and \texttt{19-XX}. Its license notice states that, unless otherwise noted, the report content is licensed under CC BY SA 4.0.
\end{itemize}
\end{document}
